% Abstract
\chapter*{Abstract}
\addcontentsline{toc}{chapter}{Abstract}

This monograph presents a formal reconstruction of the Indus (Harappan) script as an operational symbolic system implementing finite mathematics, relational logic, and state-based validation. Rather than attempting phonetic decipherment, this work analyzes the script as a formal system of operational notation used for administrative and institutional validation.

The reconstruction demonstrates that the Indus script encodes:
\begin{itemize}
    \item \textbf{Finite arithmetic}: Cardinal operations and counting systems
    \item \textbf{Relational logic}: State transitions and logical operators
    \item \textbf{Institutional validation}: Administrative control and verification mechanisms
    \item \textbf{Formal equivalents}: Structures isomorphic to finite automata and modern computational systems
\end{itemize}

Through systematic analysis of sign frequencies, distributional patterns, and combinatorial constraints, we establish that the Indus script functions as a domain-specific operational language rather than a general-purpose phonetic writing system. The mathematical formalism developed here provides a rigorous framework for understanding the script's purpose and limitations.

\textbf{Key contributions:}
\begin{enumerate}
    \item A complete operational taxonomy of Indus sign functions
    \item Formal proofs of the system's mathematical scope and constraints
    \item Cross-cultural comparison with analogous symbolic systems
    \item Computational models demonstrating functional equivalence
\end{enumerate}

This work is intended for researchers in computational archaeology, ancient mathematics, writing systems, and formal methods applied to historical problems.
